\documentclass[aspectratio=169, 11pt]{beamer}
% ============================================================
% Preambles/header.tex
% MN52233 — Causal Identification Strategies
% University of Bath, School of Management
% ============================================================
% Shared preamble for all lecture slide decks.
% Usage in each .tex file: % ============================================================
% Preambles/header.tex
% MN52233 — Causal Identification Strategies
% University of Bath, School of Management
% ============================================================
% Shared preamble for all lecture slide decks.
% Usage in each .tex file: % ============================================================
% Preambles/header.tex
% MN52233 — Causal Identification Strategies
% University of Bath, School of Management
% ============================================================
% Shared preamble for all lecture slide decks.
% Usage in each .tex file: \input{../Preambles/header.tex}
% ============================================================

% --- Core packages ---
\usepackage{amsmath,amssymb,mathtools}
\usepackage{xcolor}
\usepackage{graphicx}
\usepackage{booktabs}
\usepackage{tikz}
\usetikzlibrary{arrows.meta,positioning,shapes.geometric,calc,decorations.pathreplacing}
\usepackage{tcolorbox}
\tcbuselibrary{skins}
\usepackage{natbib}
\bibpunct{(}{)}{;}{a}{,}{,}
\usepackage{enumerate}
\usepackage{ragged2e}
\usepackage{microtype}

% ============================================================
% COLOUR DEFINITIONS
% ============================================================

% Primary palette
\definecolor{bathnav}{HTML}{1F3A5F}      % Dark navy — primary
\definecolor{bathacc}{HTML}{3A7EBF}      % Medium blue — secondary
\definecolor{bathlight}{HTML}{EDF3FB}    % Very light blue — box backgrounds
\definecolor{nearblack}{HTML}{1A1A1A}    % Near-black body text

% Okabe-Ito colorblind-safe palette (mandatory for all figures)
\definecolor{okabeblue}{HTML}{0072B2}
\definecolor{okabeorange}{HTML}{E69F00}
\definecolor{okabegreen}{HTML}{009E73}
\definecolor{okabesky}{HTML}{56B4E9}
\definecolor{okabevermillion}{HTML}{D55E00}
\definecolor{okabepink}{HTML}{CC79A7}
\definecolor{okabeyellow}{HTML}{F0E442}

% ============================================================
% BEAMER THEME CUSTOMISATION
% ============================================================

\usetheme{default}

% --- Text & structure colours ---
\setbeamercolor{normal text}{fg=nearblack, bg=white}
\setbeamercolor{structure}{fg=bathnav}
\setbeamercolor{alerted text}{fg=okabevermillion}
\setbeamercolor{example text}{fg=okabegreen}

% --- Title page ---
\setbeamercolor{title}{fg=bathnav}
\setbeamercolor{subtitle}{fg=bathacc}
\setbeamercolor{author}{fg=nearblack}
\setbeamercolor{institute}{fg=nearblack!70}
\setbeamercolor{date}{fg=bathacc}

% --- Frame title ---
\setbeamercolor{frametitle}{bg=bathnav, fg=white}
\setbeamerfont{frametitle}{series=\bfseries, size=\normalsize}

% --- Itemize bullets ---
\setbeamercolor{item}{fg=bathnav}
\setbeamercolor{subitem}{fg=bathacc}
\setbeamertemplate{itemize item}{\small\raise0.5pt\hbox{\color{bathnav}$\bullet$}}
\setbeamertemplate{itemize subitem}{\small\raise0.5pt\hbox{\color{bathacc}$\circ$}}
\setbeamertemplate{itemize subsubitem}{\scriptsize\raise1pt\hbox{\color{bathacc}$-$}}

% --- Block colours ---
\setbeamercolor{block title}{bg=bathnav, fg=white}
\setbeamercolor{block body}{bg=bathlight, fg=nearblack}
\setbeamercolor{block title alerted}{bg=okabevermillion, fg=white}
\setbeamercolor{block body alerted}{bg=okabevermillion!10, fg=nearblack}
\setbeamercolor{block title example}{bg=okabegreen, fg=white}
\setbeamercolor{block body example}{bg=okabegreen!10, fg=nearblack}

% --- Frame title template ---
\setbeamertemplate{frametitle}{%
  \nointerlineskip%
  \begin{beamercolorbox}[wd=\paperwidth, ht=2.8ex, dp=1.2ex, leftskip=1em]{frametitle}%
    \usebeamerfont{frametitle}\insertframetitle%
  \end{beamercolorbox}%
  \nointerlineskip%
  {\color{bathacc}\rule{\paperwidth}{0.4pt}}%
  \vspace{0.3em}%
}

% --- Remove navigation symbols ---
\setbeamertemplate{navigation symbols}{}

% --- Footer ---
\setbeamercolor{footline}{fg=nearblack!60, bg=white}
\setbeamertemplate{footline}{%
  \leavevmode%
  {\color{bathnav!30}\rule{\paperwidth}{0.3pt}}%
  \hbox{%
    \begin{beamercolorbox}[wd=.30\paperwidth, ht=2.2ex, dp=1ex, leftskip=0.8em]{footline}%
      \footnotesize\color{bathacc}MN52233%
    \end{beamercolorbox}%
    \begin{beamercolorbox}[wd=.40\paperwidth, ht=2.2ex, dp=1ex, center]{footline}%
      \footnotesize\color{nearblack!60}\insertshorttitle%
    \end{beamercolorbox}%
    \begin{beamercolorbox}[wd=.30\paperwidth, ht=2.2ex, dp=1ex, rightskip=0.8em, right]{footline}%
      \footnotesize\color{nearblack!60}\insertframenumber\,/\,\inserttotalframenumber%
    \end{beamercolorbox}%
  }%
  \vskip2pt%
}

% --- Title page template ---
\setbeamertemplate{title page}{%
  \vspace{1.8em}%
  \begin{beamercolorbox}[sep=6pt, center]{title}%
    \usebeamerfont{title}\usebeamercolor[fg]{title}\inserttitle\par%
  \end{beamercolorbox}%
  \vskip0.4em%
  \begin{beamercolorbox}[sep=4pt, center]{subtitle}%
    \usebeamerfont{subtitle}\usebeamercolor[fg]{subtitle}\insertsubtitle\par%
  \end{beamercolorbox}%
  \vskip0.8em%
  {\color{bathnav!40}\rule{0.5\paperwidth}{0.4pt}}\par%
  \vskip0.8em%
  \begin{beamercolorbox}[sep=3pt, center]{author}%
    \usebeamerfont{author}\insertauthor%
  \end{beamercolorbox}%
  \begin{beamercolorbox}[sep=2pt, center]{institute}%
    \usebeamerfont{institute}\insertinstitute%
  \end{beamercolorbox}%
  \vskip0.4em%
  \begin{beamercolorbox}[sep=2pt, center]{date}%
    \usebeamerfont{date}\insertdate%
  \end{beamercolorbox}%
}

% ============================================================
% SECTION TRANSITION SLIDES
% ============================================================

% Usage: \sectionslide{Section Title}{Optional subtitle or question}
\newcommand{\sectionslide}[2]{%
  \begingroup%
  \setbeamercolor{background canvas}{bg=bathnav}%
  \begin{frame}[plain, noframenumbering]%
    \vfill%
    \centering%
    {\Large\bfseries\textcolor{white}{#1}}\par%
    \vskip0.8em%
    \ifx\relax#2\relax\else%
      {\normalsize\textcolor{okabesky}{#2}}\par%
    \fi%
    \vfill%
  \end{frame}%
  \endgroup%
}

% ============================================================
% TCOLORBOX CUSTOM ENVIRONMENTS
% ============================================================

\tcbset{
  beamerbase/.style={
    enhanced,
    boxrule=0.4pt,
    arc=2pt,
    left=6pt, right=6pt, top=3pt, bottom=4pt,
    fonttitle=\bfseries\small,
    attach boxed title to top left={yshift=-0.1mm},
    boxed title style={
      arc=2pt,
      boxrule=0pt,
      left=6pt, right=6pt, top=3pt, bottom=3pt,
    },
    before upper={\setlength\parindent{0pt}\small},
  }
}

% Key point box — dark navy title bar, white title text, light blue body
\newtcolorbox{keybox}[1][Key Point]{
  beamerbase,
  colback=bathnav!8,
  colframe=bathnav,
  colbacktitle=bathnav,
  coltitle=white,
  title={#1},
}

% Paper example box — dark orange title bar, white title text, light orange body
\newtcolorbox{examplebox}[1][Example]{
  beamerbase,
  colback=okabeorange!12,
  colframe=okabeorange!80!black,
  colbacktitle=okabeorange!85!black,
  coltitle=white,
  title={#1},
}

% Warning / caveat box — vermillion title bar, white title text, light red body
\newtcolorbox{warningbox}[1][Caution]{
  beamerbase,
  colback=okabevermillion!10,
  colframe=okabevermillion,
  colbacktitle=okabevermillion,
  coltitle=white,
  title={#1},
}

% Intuition box — green title bar, white title text, light green body
\newtcolorbox{intuitionbox}[1][Intuition]{
  beamerbase,
  colback=okabegreen!10,
  colframe=okabegreen,
  colbacktitle=okabegreen,
  coltitle=white,
  title={#1},
}

% ============================================================
% MATH MACROS
% ============================================================

\newcommand{\E}[1]{\mathbb{E}\!\left[#1\right]}
\newcommand{\Econd}[2]{\mathbb{E}\!\left[#1\,\middle|\,#2\right]}
\newcommand{\Yi}[1]{Y_{i}(#1)}
\newcommand{\Y}[1]{Y(#1)}
\newcommand{\ATE}{\mathrm{ATE}}
\newcommand{\ATT}{\mathrm{ATT}}
\newcommand{\ITT}{\mathrm{ITT}}
\newcommand{\LATE}{\mathrm{LATE}}
\newcommand{\indep}{\perp\!\!\!\perp}
\newcommand{\muted}[1]{{\color{nearblack!50}#1}}
\newcommand{\highlight}[1]{{\color{bathnav}\textbf{#1}}}
\newcommand{\source}[1]{\vspace{0.3em}\par{\scriptsize\color{nearblack!60}\textit{#1}}}

% ============================================================
% BIBLIOGRAPHY STYLE
% ============================================================

\bibliographystyle{plainnat}

% ============================================================

% --- Core packages ---
\usepackage{amsmath,amssymb,mathtools}
\usepackage{xcolor}
\usepackage{graphicx}
\usepackage{booktabs}
\usepackage{tikz}
\usetikzlibrary{arrows.meta,positioning,shapes.geometric,calc,decorations.pathreplacing}
\usepackage{tcolorbox}
\tcbuselibrary{skins}
\usepackage{natbib}
\bibpunct{(}{)}{;}{a}{,}{,}
\usepackage{enumerate}
\usepackage{ragged2e}
\usepackage{microtype}

% ============================================================
% COLOUR DEFINITIONS
% ============================================================

% Primary palette
\definecolor{bathnav}{HTML}{1F3A5F}      % Dark navy — primary
\definecolor{bathacc}{HTML}{3A7EBF}      % Medium blue — secondary
\definecolor{bathlight}{HTML}{EDF3FB}    % Very light blue — box backgrounds
\definecolor{nearblack}{HTML}{1A1A1A}    % Near-black body text

% Okabe-Ito colorblind-safe palette (mandatory for all figures)
\definecolor{okabeblue}{HTML}{0072B2}
\definecolor{okabeorange}{HTML}{E69F00}
\definecolor{okabegreen}{HTML}{009E73}
\definecolor{okabesky}{HTML}{56B4E9}
\definecolor{okabevermillion}{HTML}{D55E00}
\definecolor{okabepink}{HTML}{CC79A7}
\definecolor{okabeyellow}{HTML}{F0E442}

% ============================================================
% BEAMER THEME CUSTOMISATION
% ============================================================

\usetheme{default}

% --- Text & structure colours ---
\setbeamercolor{normal text}{fg=nearblack, bg=white}
\setbeamercolor{structure}{fg=bathnav}
\setbeamercolor{alerted text}{fg=okabevermillion}
\setbeamercolor{example text}{fg=okabegreen}

% --- Title page ---
\setbeamercolor{title}{fg=bathnav}
\setbeamercolor{subtitle}{fg=bathacc}
\setbeamercolor{author}{fg=nearblack}
\setbeamercolor{institute}{fg=nearblack!70}
\setbeamercolor{date}{fg=bathacc}

% --- Frame title ---
\setbeamercolor{frametitle}{bg=bathnav, fg=white}
\setbeamerfont{frametitle}{series=\bfseries, size=\normalsize}

% --- Itemize bullets ---
\setbeamercolor{item}{fg=bathnav}
\setbeamercolor{subitem}{fg=bathacc}
\setbeamertemplate{itemize item}{\small\raise0.5pt\hbox{\color{bathnav}$\bullet$}}
\setbeamertemplate{itemize subitem}{\small\raise0.5pt\hbox{\color{bathacc}$\circ$}}
\setbeamertemplate{itemize subsubitem}{\scriptsize\raise1pt\hbox{\color{bathacc}$-$}}

% --- Block colours ---
\setbeamercolor{block title}{bg=bathnav, fg=white}
\setbeamercolor{block body}{bg=bathlight, fg=nearblack}
\setbeamercolor{block title alerted}{bg=okabevermillion, fg=white}
\setbeamercolor{block body alerted}{bg=okabevermillion!10, fg=nearblack}
\setbeamercolor{block title example}{bg=okabegreen, fg=white}
\setbeamercolor{block body example}{bg=okabegreen!10, fg=nearblack}

% --- Frame title template ---
\setbeamertemplate{frametitle}{%
  \nointerlineskip%
  \begin{beamercolorbox}[wd=\paperwidth, ht=2.8ex, dp=1.2ex, leftskip=1em]{frametitle}%
    \usebeamerfont{frametitle}\insertframetitle%
  \end{beamercolorbox}%
  \nointerlineskip%
  {\color{bathacc}\rule{\paperwidth}{0.4pt}}%
  \vspace{0.3em}%
}

% --- Remove navigation symbols ---
\setbeamertemplate{navigation symbols}{}

% --- Footer ---
\setbeamercolor{footline}{fg=nearblack!60, bg=white}
\setbeamertemplate{footline}{%
  \leavevmode%
  {\color{bathnav!30}\rule{\paperwidth}{0.3pt}}%
  \hbox{%
    \begin{beamercolorbox}[wd=.30\paperwidth, ht=2.2ex, dp=1ex, leftskip=0.8em]{footline}%
      \footnotesize\color{bathacc}MN52233%
    \end{beamercolorbox}%
    \begin{beamercolorbox}[wd=.40\paperwidth, ht=2.2ex, dp=1ex, center]{footline}%
      \footnotesize\color{nearblack!60}\insertshorttitle%
    \end{beamercolorbox}%
    \begin{beamercolorbox}[wd=.30\paperwidth, ht=2.2ex, dp=1ex, rightskip=0.8em, right]{footline}%
      \footnotesize\color{nearblack!60}\insertframenumber\,/\,\inserttotalframenumber%
    \end{beamercolorbox}%
  }%
  \vskip2pt%
}

% --- Title page template ---
\setbeamertemplate{title page}{%
  \vspace{1.8em}%
  \begin{beamercolorbox}[sep=6pt, center]{title}%
    \usebeamerfont{title}\usebeamercolor[fg]{title}\inserttitle\par%
  \end{beamercolorbox}%
  \vskip0.4em%
  \begin{beamercolorbox}[sep=4pt, center]{subtitle}%
    \usebeamerfont{subtitle}\usebeamercolor[fg]{subtitle}\insertsubtitle\par%
  \end{beamercolorbox}%
  \vskip0.8em%
  {\color{bathnav!40}\rule{0.5\paperwidth}{0.4pt}}\par%
  \vskip0.8em%
  \begin{beamercolorbox}[sep=3pt, center]{author}%
    \usebeamerfont{author}\insertauthor%
  \end{beamercolorbox}%
  \begin{beamercolorbox}[sep=2pt, center]{institute}%
    \usebeamerfont{institute}\insertinstitute%
  \end{beamercolorbox}%
  \vskip0.4em%
  \begin{beamercolorbox}[sep=2pt, center]{date}%
    \usebeamerfont{date}\insertdate%
  \end{beamercolorbox}%
}

% ============================================================
% SECTION TRANSITION SLIDES
% ============================================================

% Usage: \sectionslide{Section Title}{Optional subtitle or question}
\newcommand{\sectionslide}[2]{%
  \begingroup%
  \setbeamercolor{background canvas}{bg=bathnav}%
  \begin{frame}[plain, noframenumbering]%
    \vfill%
    \centering%
    {\Large\bfseries\textcolor{white}{#1}}\par%
    \vskip0.8em%
    \ifx\relax#2\relax\else%
      {\normalsize\textcolor{okabesky}{#2}}\par%
    \fi%
    \vfill%
  \end{frame}%
  \endgroup%
}

% ============================================================
% TCOLORBOX CUSTOM ENVIRONMENTS
% ============================================================

\tcbset{
  beamerbase/.style={
    enhanced,
    boxrule=0.4pt,
    arc=2pt,
    left=6pt, right=6pt, top=3pt, bottom=4pt,
    fonttitle=\bfseries\small,
    attach boxed title to top left={yshift=-0.1mm},
    boxed title style={
      arc=2pt,
      boxrule=0pt,
      left=6pt, right=6pt, top=3pt, bottom=3pt,
    },
    before upper={\setlength\parindent{0pt}\small},
  }
}

% Key point box — dark navy title bar, white title text, light blue body
\newtcolorbox{keybox}[1][Key Point]{
  beamerbase,
  colback=bathnav!8,
  colframe=bathnav,
  colbacktitle=bathnav,
  coltitle=white,
  title={#1},
}

% Paper example box — dark orange title bar, white title text, light orange body
\newtcolorbox{examplebox}[1][Example]{
  beamerbase,
  colback=okabeorange!12,
  colframe=okabeorange!80!black,
  colbacktitle=okabeorange!85!black,
  coltitle=white,
  title={#1},
}

% Warning / caveat box — vermillion title bar, white title text, light red body
\newtcolorbox{warningbox}[1][Caution]{
  beamerbase,
  colback=okabevermillion!10,
  colframe=okabevermillion,
  colbacktitle=okabevermillion,
  coltitle=white,
  title={#1},
}

% Intuition box — green title bar, white title text, light green body
\newtcolorbox{intuitionbox}[1][Intuition]{
  beamerbase,
  colback=okabegreen!10,
  colframe=okabegreen,
  colbacktitle=okabegreen,
  coltitle=white,
  title={#1},
}

% ============================================================
% MATH MACROS
% ============================================================

\newcommand{\E}[1]{\mathbb{E}\!\left[#1\right]}
\newcommand{\Econd}[2]{\mathbb{E}\!\left[#1\,\middle|\,#2\right]}
\newcommand{\Yi}[1]{Y_{i}(#1)}
\newcommand{\Y}[1]{Y(#1)}
\newcommand{\ATE}{\mathrm{ATE}}
\newcommand{\ATT}{\mathrm{ATT}}
\newcommand{\ITT}{\mathrm{ITT}}
\newcommand{\LATE}{\mathrm{LATE}}
\newcommand{\indep}{\perp\!\!\!\perp}
\newcommand{\muted}[1]{{\color{nearblack!50}#1}}
\newcommand{\highlight}[1]{{\color{bathnav}\textbf{#1}}}
\newcommand{\source}[1]{\vspace{0.3em}\par{\scriptsize\color{nearblack!60}\textit{#1}}}

% ============================================================
% BIBLIOGRAPHY STYLE
% ============================================================

\bibliographystyle{plainnat}

% ============================================================

% --- Core packages ---
\usepackage{amsmath,amssymb,mathtools}
\usepackage{xcolor}
\usepackage{graphicx}
\usepackage{booktabs}
\usepackage{tikz}
\usetikzlibrary{arrows.meta,positioning,shapes.geometric,calc,decorations.pathreplacing}
\usepackage{tcolorbox}
\tcbuselibrary{skins}
\usepackage{natbib}
\bibpunct{(}{)}{;}{a}{,}{,}
\usepackage{enumerate}
\usepackage{ragged2e}
\usepackage{microtype}

% ============================================================
% COLOUR DEFINITIONS
% ============================================================

% Primary palette
\definecolor{bathnav}{HTML}{1F3A5F}      % Dark navy — primary
\definecolor{bathacc}{HTML}{3A7EBF}      % Medium blue — secondary
\definecolor{bathlight}{HTML}{EDF3FB}    % Very light blue — box backgrounds
\definecolor{nearblack}{HTML}{1A1A1A}    % Near-black body text

% Okabe-Ito colorblind-safe palette (mandatory for all figures)
\definecolor{okabeblue}{HTML}{0072B2}
\definecolor{okabeorange}{HTML}{E69F00}
\definecolor{okabegreen}{HTML}{009E73}
\definecolor{okabesky}{HTML}{56B4E9}
\definecolor{okabevermillion}{HTML}{D55E00}
\definecolor{okabepink}{HTML}{CC79A7}
\definecolor{okabeyellow}{HTML}{F0E442}

% ============================================================
% BEAMER THEME CUSTOMISATION
% ============================================================

\usetheme{default}

% --- Text & structure colours ---
\setbeamercolor{normal text}{fg=nearblack, bg=white}
\setbeamercolor{structure}{fg=bathnav}
\setbeamercolor{alerted text}{fg=okabevermillion}
\setbeamercolor{example text}{fg=okabegreen}

% --- Title page ---
\setbeamercolor{title}{fg=bathnav}
\setbeamercolor{subtitle}{fg=bathacc}
\setbeamercolor{author}{fg=nearblack}
\setbeamercolor{institute}{fg=nearblack!70}
\setbeamercolor{date}{fg=bathacc}

% --- Frame title ---
\setbeamercolor{frametitle}{bg=bathnav, fg=white}
\setbeamerfont{frametitle}{series=\bfseries, size=\normalsize}

% --- Itemize bullets ---
\setbeamercolor{item}{fg=bathnav}
\setbeamercolor{subitem}{fg=bathacc}
\setbeamertemplate{itemize item}{\small\raise0.5pt\hbox{\color{bathnav}$\bullet$}}
\setbeamertemplate{itemize subitem}{\small\raise0.5pt\hbox{\color{bathacc}$\circ$}}
\setbeamertemplate{itemize subsubitem}{\scriptsize\raise1pt\hbox{\color{bathacc}$-$}}

% --- Block colours ---
\setbeamercolor{block title}{bg=bathnav, fg=white}
\setbeamercolor{block body}{bg=bathlight, fg=nearblack}
\setbeamercolor{block title alerted}{bg=okabevermillion, fg=white}
\setbeamercolor{block body alerted}{bg=okabevermillion!10, fg=nearblack}
\setbeamercolor{block title example}{bg=okabegreen, fg=white}
\setbeamercolor{block body example}{bg=okabegreen!10, fg=nearblack}

% --- Frame title template ---
\setbeamertemplate{frametitle}{%
  \nointerlineskip%
  \begin{beamercolorbox}[wd=\paperwidth, ht=2.8ex, dp=1.2ex, leftskip=1em]{frametitle}%
    \usebeamerfont{frametitle}\insertframetitle%
  \end{beamercolorbox}%
  \nointerlineskip%
  {\color{bathacc}\rule{\paperwidth}{0.4pt}}%
  \vspace{0.3em}%
}

% --- Remove navigation symbols ---
\setbeamertemplate{navigation symbols}{}

% --- Footer ---
\setbeamercolor{footline}{fg=nearblack!60, bg=white}
\setbeamertemplate{footline}{%
  \leavevmode%
  {\color{bathnav!30}\rule{\paperwidth}{0.3pt}}%
  \hbox{%
    \begin{beamercolorbox}[wd=.30\paperwidth, ht=2.2ex, dp=1ex, leftskip=0.8em]{footline}%
      \footnotesize\color{bathacc}MN52233%
    \end{beamercolorbox}%
    \begin{beamercolorbox}[wd=.40\paperwidth, ht=2.2ex, dp=1ex, center]{footline}%
      \footnotesize\color{nearblack!60}\insertshorttitle%
    \end{beamercolorbox}%
    \begin{beamercolorbox}[wd=.30\paperwidth, ht=2.2ex, dp=1ex, rightskip=0.8em, right]{footline}%
      \footnotesize\color{nearblack!60}\insertframenumber\,/\,\inserttotalframenumber%
    \end{beamercolorbox}%
  }%
  \vskip2pt%
}

% --- Title page template ---
\setbeamertemplate{title page}{%
  \vspace{1.8em}%
  \begin{beamercolorbox}[sep=6pt, center]{title}%
    \usebeamerfont{title}\usebeamercolor[fg]{title}\inserttitle\par%
  \end{beamercolorbox}%
  \vskip0.4em%
  \begin{beamercolorbox}[sep=4pt, center]{subtitle}%
    \usebeamerfont{subtitle}\usebeamercolor[fg]{subtitle}\insertsubtitle\par%
  \end{beamercolorbox}%
  \vskip0.8em%
  {\color{bathnav!40}\rule{0.5\paperwidth}{0.4pt}}\par%
  \vskip0.8em%
  \begin{beamercolorbox}[sep=3pt, center]{author}%
    \usebeamerfont{author}\insertauthor%
  \end{beamercolorbox}%
  \begin{beamercolorbox}[sep=2pt, center]{institute}%
    \usebeamerfont{institute}\insertinstitute%
  \end{beamercolorbox}%
  \vskip0.4em%
  \begin{beamercolorbox}[sep=2pt, center]{date}%
    \usebeamerfont{date}\insertdate%
  \end{beamercolorbox}%
}

% ============================================================
% SECTION TRANSITION SLIDES
% ============================================================

% Usage: \sectionslide{Section Title}{Optional subtitle or question}
\newcommand{\sectionslide}[2]{%
  \begingroup%
  \setbeamercolor{background canvas}{bg=bathnav}%
  \begin{frame}[plain, noframenumbering]%
    \vfill%
    \centering%
    {\Large\bfseries\textcolor{white}{#1}}\par%
    \vskip0.8em%
    \ifx\relax#2\relax\else%
      {\normalsize\textcolor{okabesky}{#2}}\par%
    \fi%
    \vfill%
  \end{frame}%
  \endgroup%
}

% ============================================================
% TCOLORBOX CUSTOM ENVIRONMENTS
% ============================================================

\tcbset{
  beamerbase/.style={
    enhanced,
    boxrule=0.4pt,
    arc=2pt,
    left=6pt, right=6pt, top=3pt, bottom=4pt,
    fonttitle=\bfseries\small,
    attach boxed title to top left={yshift=-0.1mm},
    boxed title style={
      arc=2pt,
      boxrule=0pt,
      left=6pt, right=6pt, top=3pt, bottom=3pt,
    },
    before upper={\setlength\parindent{0pt}\small},
  }
}

% Key point box — dark navy title bar, white title text, light blue body
\newtcolorbox{keybox}[1][Key Point]{
  beamerbase,
  colback=bathnav!8,
  colframe=bathnav,
  colbacktitle=bathnav,
  coltitle=white,
  title={#1},
}

% Paper example box — dark orange title bar, white title text, light orange body
\newtcolorbox{examplebox}[1][Example]{
  beamerbase,
  colback=okabeorange!12,
  colframe=okabeorange!80!black,
  colbacktitle=okabeorange!85!black,
  coltitle=white,
  title={#1},
}

% Warning / caveat box — vermillion title bar, white title text, light red body
\newtcolorbox{warningbox}[1][Caution]{
  beamerbase,
  colback=okabevermillion!10,
  colframe=okabevermillion,
  colbacktitle=okabevermillion,
  coltitle=white,
  title={#1},
}

% Intuition box — green title bar, white title text, light green body
\newtcolorbox{intuitionbox}[1][Intuition]{
  beamerbase,
  colback=okabegreen!10,
  colframe=okabegreen,
  colbacktitle=okabegreen,
  coltitle=white,
  title={#1},
}

% ============================================================
% MATH MACROS
% ============================================================

\newcommand{\E}[1]{\mathbb{E}\!\left[#1\right]}
\newcommand{\Econd}[2]{\mathbb{E}\!\left[#1\,\middle|\,#2\right]}
\newcommand{\Yi}[1]{Y_{i}(#1)}
\newcommand{\Y}[1]{Y(#1)}
\newcommand{\ATE}{\mathrm{ATE}}
\newcommand{\ATT}{\mathrm{ATT}}
\newcommand{\ITT}{\mathrm{ITT}}
\newcommand{\LATE}{\mathrm{LATE}}
\newcommand{\indep}{\perp\!\!\!\perp}
\newcommand{\muted}[1]{{\color{nearblack!50}#1}}
\newcommand{\highlight}[1]{{\color{bathnav}\textbf{#1}}}
\newcommand{\source}[1]{\vspace{0.3em}\par{\scriptsize\color{nearblack!60}\textit{#1}}}

% ============================================================
% BIBLIOGRAPHY STYLE
% ============================================================

\bibliographystyle{plainnat}


% ============================================================
% METADATA
% ============================================================
\title[D1A: Lab Experiments \& RCTs]{Lab Experiments and Randomized Controlled Trials}
\subtitle{Causal Identification Strategies}
\author{Na Zou}
\institute{University of Bath \\ School of Management}
\date{MN52233 \quad Day 1 $\cdot$ Session A}

% ============================================================
\begin{document}
% ============================================================

% --- Title slide ---
\begin{frame}[plain, noframenumbering]
  \titlepage
\end{frame}

% ============================================================
% COURSE ROADMAP
% ============================================================

\begin{frame}[t]{Course Roadmap: Three Days, Three Strategies}
\vspace{0.5em}
\begin{center}
\footnotesize
\begin{tabular}{lll}
\toprule
\textbf{Day} & \textbf{Session A (Concepts)} & \textbf{Session B (Lab / Replication)} \\
\midrule
Day 1 & \highlight{Lab Experiments \& RCTs} & Settele (2022) — survey experiment \\
Day 2 & Regression Discontinuity \& DiD & Braghieri et al. (2022) — DiD \\
Day 3 & Instrumental Variables \& PSM & Aghion et al. (2013) — IV \\
\bottomrule
\end{tabular}
\end{center}
\vspace{0.8em}
\begin{keybox}[Central question across all three days]
Under what conditions can we give a \textbf{causal} interpretation to the
relationship between a treatment $D$ and an outcome $Y$?
\end{keybox}
\end{frame}

% --- Today's roadmap ---
\begin{frame}[t]{Today's Roadmap}
\begin{enumerate}[\quad 1.]
  \setlength\itemsep{0.5em}
  \item \textbf{Why causal inference?} — the identification problem
  \item \textbf{Potential outcomes framework} — the language of causality
  \item \textbf{Randomized Controlled Trials} — the gold standard
  \item \textbf{Lab, field, and survey experiments} — the family tree
  \item \textbf{Complications in practice} — what can go wrong
  \item \textbf{Validity} — what experiments can and cannot tell us
\end{enumerate}
\vspace{0.6em}
\muted{\small Session B (13:00--14:40): STATA replication of \citet{Settele2022_gender}}
\end{frame}

% ============================================================
\sectionslide{1. Why Causal Inference?}{The identification problem}
% ============================================================

\begin{frame}[t]{The Central Question}
\begin{itemize}
  \setlength\itemsep{0.6em}
  \item We observe that $D$ and $Y$ move together in data.
  \item Does $D$ \textbf{cause} $Y$? Or do they co-move for other reasons?
  \item This distinction has real consequences:
  \begin{itemize}
    \item Should we \emph{expand} a job training programme?
    \item Does a new drug \emph{reduce} hospitalisation?
    \item Does higher education \emph{increase} wages?
  \end{itemize}
\end{itemize}
\vspace{0.5em}
\begin{intuitionbox}
  Correlation tells us two things move together.\\
  Causation tells us \emph{why} — and whether intervening would change outcomes.
\end{intuitionbox}
\end{frame}

\begin{frame}[t]{A Striking Non-Example}
\begin{columns}[T]
\begin{column}{0.55\textwidth}
\begin{itemize}
  \setlength\itemsep{0.5em}
  \item On hot summer days, ice cream sales are high.
  \item On those same days, drowning incidents are also high.
  \item Should we \textbf{ban ice cream} to prevent drownings?
\end{itemize}
\vspace{0.5em}
\begin{warningbox}[The omitted variable]
Hot weather drives both. Ice cream does not cause drowning.
The causal link runs from \emph{temperature}, not ice cream.
\end{warningbox}
\end{column}
\begin{column}{0.42\textwidth}
\begin{center}
\begin{tikzpicture}[scale=0.82, every node/.style={font=\small}]
  \node[draw, rounded corners, fill=okabesky!20, minimum width=2cm, minimum height=0.7cm] (T) at (0,2) {Temperature};
  \node[draw, rounded corners, fill=okabeorange!20, minimum width=2cm, minimum height=0.7cm] (I) at (-1.5,0) {Ice cream sales};
  \node[draw, rounded corners, fill=okabevermillion!20, minimum width=2cm, minimum height=0.7cm] (D) at (1.5,0) {Drownings};
  \draw[->, thick, color=okabesky!80!black] (T) -- (I);
  \draw[->, thick, color=okabesky!80!black] (T) -- (D);
  \draw[<->, dashed, thick, color=nearblack!40] (I) -- node[below, font=\scriptsize, color=nearblack!60]{correlation} (D);
\end{tikzpicture}
\end{center}
\end{column}
\end{columns}
\end{frame}

\begin{frame}[t]{A Policy-Relevant Example}
\begin{itemize}
  \setlength\itemsep{0.5em}
  \item University graduates earn, on average, substantially more than non-graduates.
  \item Naive interpretation: \emph{going to university raises wages}.
  \item Is this causal?
\end{itemize}
\vspace{0.4em}
\begin{warningbox}[Selection problem]
People who attend university are different \emph{in many ways} from those who do not:\\
ability, family background, career aspirations, social networks, \ldots\\[0.3em]
The wage gap may reflect these differences, not the education itself.
\end{warningbox}
\vspace{0.3em}
\muted{\small This is the \textbf{selection bias} problem — the core challenge of causal inference.}
\end{frame}

\begin{frame}[t]{The Fundamental Problem of Causal Inference}
\begin{columns}[T]
\begin{column}{0.55\textwidth}
\begin{itemize}
  \setlength\itemsep{0.5em}
  \item For any individual $i$, we want to know:
  \[
    \tau_i = \Yi{1} - \Yi{0}
  \]
  \item $\Yi{1}$: outcome \emph{if} treated. \quad $\Yi{0}$: outcome \emph{if not} treated.
  \item \textbf{Problem:} We can only ever observe \emph{one} of these.
  \item The unobserved outcome is the \textbf{counterfactual}.
\end{itemize}
\vspace{0.4em}
\begin{keybox}[Holland (1986)]
``The fundamental problem of causal inference is that we cannot observe
the same unit in two states of the world simultaneously.''
\end{keybox}
\end{column}
\begin{column}{0.42\textwidth}
\begin{center}
\begin{tikzpicture}[scale=0.82, every node/.style={font=\small}]
  \node[draw, circle, fill=bathnav!10, minimum size=1.3cm, thick, color=bathnav] (i) at (0,0) {\small Unit $i$};
  \node[draw, rounded corners, fill=okabeblue!15, minimum width=2.2cm, minimum height=0.7cm, color=okabeblue, thick] (y1) at (3.2, 1.3) {$\Yi{1}$};
  \node[font=\scriptsize, color=okabeblue] at (3.2, 0.7) {Observed $\checkmark$};
  \node[draw, rounded corners, fill=nearblack!8, minimum width=2.2cm, minimum height=0.7cm, color=nearblack!40, thick, dashed] (y0) at (3.2, -1.3) {$\Yi{0}$};
  \node[font=\scriptsize, color=nearblack!50] at (3.2, -1.9) {Counterfactual \texttimes};
  \draw[->, thick, color=okabeblue] (i) -- (y1) node[midway, above, sloped, font=\scriptsize]{$D_i=1$};
  \draw[->, thick, dashed, color=nearblack!40] (i) -- (y0) node[midway, below, sloped, font=\scriptsize]{$D_i=0$};
\end{tikzpicture}
\end{center}
\end{column}
\end{columns}
\end{frame}

\begin{frame}[t]{Counterfactuals: The Question Behind Every Causal Claim}
\begin{itemize}
  \setlength\itemsep{0.5em}
  \item Every causal question is really a \textbf{comparison} question:
  \begin{itemize}
    \item What would this patient's health be \emph{had they not taken the drug}?
    \item What would this worker's wage be \emph{had they not gone to university}?
    \item What would this country's growth be \emph{had the policy not been implemented}?
  \end{itemize}
  \item The counterfactual is the missing half of the comparison.
  \item The challenge: we need to construct a credible stand-in for the counterfactual.
\end{itemize}
\vspace{0.4em}
\begin{intuitionbox}
Causal inference is fundamentally about finding a \textbf{valid comparison group}
— units that approximate what the treated units would have looked like \emph{absent treatment}.
\end{intuitionbox}
\end{frame}

\begin{frame}[t]{Selection Bias}
\begin{itemize}
  \setlength\itemsep{0.45em}
  \item In practice, we compare treated units to untreated units.
  \item This comparison is valid only if the groups are comparable.
  \item When treated units differ from control units in ways that also affect $Y$,
        we have \textbf{selection bias}.
\end{itemize}
\vspace{0.3em}
\begin{examplebox}[\citeauthor{Bertrand2004_emily} (\citeyear{Bertrand2004_emily}): motivation]
  \small
  Applicants with ``Black-sounding'' names get fewer callbacks than those with
  ``white-sounding'' names. Is this discrimination, or do these applicants differ
  in unobserved ways?\\[0.3em]
  Na\"ive comparison confounds \emph{name} with everything else that differs
  between applicants.
\end{examplebox}
\vspace{0.3em}
\muted{\small We will see how they solve this problem on the next slides.}
\end{frame}

\begin{frame}[t]{Why Naive OLS Comparisons Mislead}
\vspace{0.2em}
The observed difference between treated and untreated decomposes as:
\[
  \underbrace{\Econd{Y_i}{D_i=1} - \Econd{Y_i}{D_i=0}}_{\text{observed gap}}
  \;=\;
  \underbrace{\Econd{\Yi{1}}{D_i=1} - \Econd{\Yi{0}}{D_i=1}}_{\text{ATT (what we want)}}
  \;+\;
  \underbrace{\Econd{\Yi{0}}{D_i=1} - \Econd{\Yi{0}}{D_i=0}}_{\text{selection bias}}
\]
\vspace{0.3em}
\begin{itemize}
  \setlength\itemsep{0.4em}
  \item The \textbf{selection bias} term is non-zero whenever untreated units differ
        from treated units in their \emph{untreated potential outcome}.
  \item OLS with controls can reduce but rarely \emph{eliminate} selection bias.
  \item We need an identification strategy that makes selection bias zero.
\end{itemize}
\begin{keybox}
The goal of causal identification is to make the selection bias term vanish.
\end{keybox}
\end{frame}

% ============================================================
\sectionslide{2. The Potential Outcomes Framework}{The language of causal inference}
% ============================================================

\begin{frame}[t]{Two Potential Worlds}
\begin{itemize}
  \setlength\itemsep{0.5em}
  \item The \textbf{Rubin Causal Model} \citep{AngristPischke2009} gives us a precise language.
  \item For each unit $i$ and each treatment state $d \in \{0,1\}$, define:
  \[
    \Yi{d} = \text{potential outcome of unit } i \text{ under treatment } d
  \]
  \item Treatment indicator: $D_i \in \{0,1\}$
  \item What we observe:
  \[
    Y_i = D_i \,\Yi{1} + (1-D_i)\,\Yi{0}
  \]
  \item We observe $\Yi{1}$ for treated units, $\Yi{0}$ for untreated units.
  \item \textbf{Never both} for the same unit.
\end{itemize}
\end{frame}

\begin{frame}[t]{Individual Treatment Effects}
\begin{itemize}
  \setlength\itemsep{0.5em}
  \item The \textbf{individual treatment effect} for unit $i$:
  \[
    \tau_i = \Yi{1} - \Yi{0}
  \]
  \item This is what we would ideally like to know for every person.
  \item It is \textbf{never observable} — we only see one side of the coin.
  \item Treatment effects may differ across individuals:
    \begin{itemize}
      \item Some benefit a lot, some a little, some may be harmed.
    \end{itemize}
\end{itemize}
\vspace{0.4em}
\begin{intuitionbox}
This is why we focus on \emph{averages}: we cannot learn $\tau_i$,
but under the right conditions we can learn $\E{\tau_i}$.
\end{intuitionbox}
\end{frame}

\begin{frame}[t]{Population Estimands: ATE, ATT, ATC}
\vspace{0.3em}
\begin{itemize}
  \setlength\itemsep{0.5em}
  \item \textbf{Average Treatment Effect (ATE)}
  \[
    \ATE = \E{\Yi{1} - \Yi{0}} = \E{\Yi{1}} - \E{\Yi{0}}
  \]
  Effect for a \emph{randomly chosen} member of the population.
  \item \textbf{Average Treatment Effect on the Treated (ATT)}
  \[
    \ATT = \Econd{\Yi{1} - \Yi{0}}{D_i = 1}
  \]
  Effect for those who \emph{actually received} treatment.
  \item \textbf{Average Treatment Effect on the Controls (ATC)}: effect for untreated.
\end{itemize}
\vspace{0.3em}
\begin{intuitionbox}
ATT answers: ``Did this policy work for the people who got it?''\\
ATE asks: ``If we rolled it out to everyone, what would happen on average?''
\end{intuitionbox}
\end{frame}

\begin{frame}[t]{SUTVA: The Hidden Assumption}
\begin{itemize}
  \setlength\itemsep{0.5em}
  \item The potential outcomes framework rests on a strong assumption:
        the \textbf{Stable Unit Treatment Value Assumption (SUTVA)}.
  \item SUTVA has two parts:
\end{itemize}
\vspace{0.4em}
\begin{columns}[T]
\begin{column}{0.48\textwidth}
  \begin{keybox}[Part 1: No interference]
    Unit $i$'s potential outcomes depend only on \emph{their own} treatment status,
    not on the treatment of others.
  \end{keybox}
\end{column}
\begin{column}{0.48\textwidth}
  \begin{keybox}[Part 2: No multiple versions]
    There is a single, well-defined version of the treatment.
    $D_i = 1$ means the same thing for all units.
  \end{keybox}
\end{column}
\end{columns}
\vspace{0.5em}
\begin{warningbox}
SUTVA is often \emph{assumed} without being stated. When it fails,
the notation $\Yi{d}$ breaks down and standard estimators are biased.
\end{warningbox}
\end{frame}

\begin{frame}[t]{When SUTVA Fails: Vaccine Trials}
\begin{itemize}
  \setlength\itemsep{0.5em}
  \item Suppose we randomly vaccinate 50\% of a community and measure infection rates.
  \item SUTVA Part 1 (no interference) is violated here:
    \begin{itemize}
      \item Vaccinated individuals are less likely to \emph{transmit} the disease.
      \item Unvaccinated individuals benefit indirectly (\textbf{herd immunity}).
      \item $\Yi{0}$ for an unvaccinated person depends on how many others are vaccinated.
    \end{itemize}
\end{itemize}
\vspace{0.4em}
\begin{examplebox}[Consequence]
If herd immunity is present, comparing infected rates in vaccinated vs.\ unvaccinated
\emph{underestimates} the true benefit of vaccinating everyone.\\[0.3em]
The ``treatment effect'' we estimate is not the one we want.
\end{examplebox}
\end{frame}

\begin{frame}[t]{Selection Bias Revisited: The Decomposition}
\vspace{0.2em}
Recall:
\[
  \underbrace{\Econd{Y_i}{D_i=1} - \Econd{Y_i}{D_i=0}}_{\text{observed gap}}
  \;=\;
  \underbrace{\ATT}_{\text{causal effect}}
  \;+\;
  \underbrace{\Econd{\Yi{0}}{D_i=1} - \Econd{\Yi{0}}{D_i=0}}_{\text{selection bias}}
\]
\vspace{0.3em}
\begin{itemize}
  \setlength\itemsep{0.4em}
  \item Selection bias is non-zero when treated units would have had different
        \emph{untreated} outcomes than untreated units.
  \item This is the \emph{baseline difference} between the two groups.
\end{itemize}
\vspace{0.3em}
\begin{keybox}[The key insight]
If $D_i \indep (\Yi{0}, \Yi{1})$, then selection bias $= 0$
and the observed gap identifies the \ATE.\\
Randomization delivers this independence.
\end{keybox}
\end{frame}

% ============================================================
\sectionslide{3. Randomized Controlled Trials}{Randomization as the gold standard}
% ============================================================

\begin{frame}[t]{What Randomization Delivers}
\begin{itemize}
  \setlength\itemsep{0.5em}
  \item In a randomized experiment, treatment $D_i$ is assigned \emph{by the researcher},
        independently of everything else.
  \item This gives us the \textbf{independence condition}:
  \[
    D_i \indep \bigl(\Yi{0},\,\Yi{1}\bigr)
  \]
  \item Under independence:
  \[
    \Econd{\Yi{0}}{D_i=1} = \Econd{\Yi{0}}{D_i=0} = \E{\Yi{0}}
  \]
  \item Selection bias vanishes: the observed gap equals the \ATE.
\end{itemize}
\vspace{0.3em}
\begin{intuitionbox}
Randomization makes treated and control groups \textbf{exchangeable}
— they are, in expectation, identical on all characteristics.
\end{intuitionbox}
\end{frame}

\begin{frame}[t]{Balance: What Randomization Looks Like in Data}
\begin{itemize}
  \setlength\itemsep{0.5em}
  \item With successful randomization, treated and control groups should be
        \textbf{similar on all pre-treatment characteristics}.
  \item This is the \textbf{randomization check} (or balance test):
    \begin{itemize}
      \item Run regressions of each pre-treatment covariate on $D_i$.
      \item We expect no significant differences.
      \item Report this \emph{always} — failure is a red flag.
    \end{itemize}
  \item Balance does not prove randomization worked, but imbalance is strong evidence it did not.
\end{itemize}
\vspace{0.3em}
\muted{\small A well-written paper always includes a ``Table 1'' showing balance across groups.}
\end{frame}

\begin{frame}[t]{Field Experiment: \citeauthor{Bertrand2004_emily} (\citeyear{Bertrand2004_emily})}
\begin{examplebox}[Design: Are Emily and Greg more employable than Lakisha and Jamal?]
  \small
  \begin{itemize}
    \item \textbf{Setting:} 1,300 newspaper job ads in Chicago and Boston
    \item \textbf{Treatment:} Resumes randomly assigned either a ``white-sounding''
          name (Emily, Greg, \ldots) or a ``Black-sounding'' name
          (Lakisha, Jamal, \ldots)
    \item \textbf{Outcome:} Callback rate for interview
    \item \textbf{Everything else} on the resume held constant or randomized
    \item \textbf{Balance:} By construction — same resumes, only name varies
  \end{itemize}
\end{examplebox}
\source{\citet{Bertrand2004_emily}}
\end{frame}

\begin{frame}[t]{\citeauthor{Bertrand2004_emily}: Main Finding}
\begin{itemize}
  \setlength\itemsep{0.5em}
  \item Resumes with white-sounding names received \textbf{50\% more callbacks}
        than identical resumes with Black-sounding names
        (9.65\% vs.\ 6.45\%).
  \item The gap is present across:
    \begin{itemize}
      \item occupation types (administrative, sales, clerical)
      \item industries
      \item employer size
    \end{itemize}
  \item Higher-quality resumes help white-sounding names more than
        Black-sounding names.
\end{itemize}
\vspace{0.3em}
\begin{keybox}[Why we believe this is causal]
Randomization ensures the \emph{only} systematic difference between the two groups
is the name on the resume. Any callback gap must reflect differential treatment
by name — i.e., discrimination.
\end{keybox}
\source{\citet{Bertrand2004_emily}}
\end{frame}

\begin{frame}[t]{Field Experiment: \citeauthor{Jessoe2014_energy} (\citeyear{Jessoe2014_energy})}
\begin{examplebox}[Design: Knowledge is (Less) Power]
  \small
  \begin{itemize}
    \item \textbf{Setting:} Residential households in Sacramento, California
    \item \textbf{Treatment:} Random assignment of ``Oracles'' — in-home displays
          showing real-time electricity usage and costs
    \item \textbf{Control:} Identical households without in-home displays
    \item \textbf{Outcome:} Electricity consumption during peak pricing periods
    \item \textbf{Key feature:} High-price days announced in advance;
          treated households have information, control households do not
  \end{itemize}
\end{examplebox}
\source{\citet{Jessoe2014_energy}}
\end{frame}

\begin{frame}[t]{\citeauthor{Jessoe2014_energy}: Main Finding}
\begin{itemize}
  \setlength\itemsep{0.5em}
  \item Treated households reduced electricity consumption by \textbf{8\% more}
        than control households during peak-price hours.
  \item Without information, demand for electricity is highly inelastic
        (people do not adjust much to price signals).
  \item With information, demand becomes significantly more elastic.
\end{itemize}
\vspace{0.3em}
\begin{intuitionbox}
The experiment isolates \textbf{information} as the mechanism.\\[0.3em]
Without randomization, households who choose to install monitors
might already be more energy-conscious — creating selection bias.
Randomization breaks that link.
\end{intuitionbox}
\source{\citet{Jessoe2014_energy}}
\end{frame}

\begin{frame}[t]{Key Design Choices in an RCT}
\begin{itemize}
  \setlength\itemsep{0.5em}
  \item \textbf{Unit of randomization:} individual, household, firm, village, school?
    \begin{itemize}
      \item Must be chosen to minimise spillovers (SUTVA concerns)
    \end{itemize}
  \item \textbf{Sample size:} determined by \emph{power calculations}
    \begin{itemize}
      \item How large an effect are we looking for? At what significance level?
    \end{itemize}
  \item \textbf{Stratified randomisation:} assign within groups (gender, region, \ldots)
    to ensure balance on key dimensions
  \item \textbf{Compliance:} do assigned units actually take up treatment?
    \begin{itemize}
      \item Non-compliance introduces complications (Section 5)
    \end{itemize}
\end{itemize}
\vspace{0.3em}
\muted{\small Good experimental design resolves most of these before data collection begins.}
\end{frame}

% ============================================================
\sectionslide{4. Lab, Field, and Survey Experiments}{The experiment family tree}
% ============================================================

\begin{frame}[t]{A Taxonomy of Experiments}
\vspace{0.3em}
\begin{center}
\begin{tikzpicture}[scale=0.78, every node/.style={font=\small}]
  \node[draw, rounded corners, fill=bathnav!15, minimum width=2.5cm, minimum height=0.9cm, thick, color=bathnav]
    (lab) at (-4,0) {\textbf{Lab}};
  \node[draw, rounded corners, fill=bathacc!15, minimum width=2.5cm, minimum height=0.9cm, thick, color=bathacc]
    (field) at (0,0) {\textbf{Field}};
  \node[draw, rounded corners, fill=okabeorange!20, minimum width=2.5cm, minimum height=0.9cm, thick, color=okabeorange!80!black]
    (survey) at (4,0) {\textbf{Survey}};

  \node[font=\scriptsize, text width=2.8cm, align=center] at (-4,-1.2) {High control\\Stylized setting\\Student subjects};
  \node[font=\scriptsize, text width=2.8cm, align=center] at (0,-1.2) {Real world\\Real subjects\\Real stakes};
  \node[font=\scriptsize, text width=2.8cm, align=center] at (4,-1.2) {Beliefs \& prefs\\Hypothetical\\Scalable};

  \draw[<->, thick, color=nearblack!50] (-2.8,0) -- (-1.2,0) node[midway, above, font=\scriptsize]{more realistic};
  \draw[<->, thick, color=nearblack!50] (1.2,0) -- (2.8,0) node[midway, above, font=\scriptsize]{cheaper/faster};
\end{tikzpicture}
\end{center}
\vspace{0.4em}
\begin{keybox}
Choosing an experiment type means trading off \textbf{internal validity}
(how clean is the causal inference?) against \textbf{external validity}
(how generalisable are the results?).
\end{keybox}
\end{frame}

\begin{frame}[t]{Lab Experiments}
\begin{columns}[T]
\begin{column}{0.55\textwidth}
\begin{itemize}
  \setlength\itemsep{0.4em}
  \item Conducted in a \textbf{controlled laboratory} setting.
  \item Researcher controls everything:
    \begin{itemize}
      \item Who participates (often students)
      \item The exact stimuli and payoffs
      \item The environment and information
    \end{itemize}
  \item Randomisation is easy and clean.
\end{itemize}
\vspace{0.4em}
\begin{columns}[T]
\begin{column}{0.48\textwidth}
  \begin{keybox}[Strength]
    Precise isolation of mechanisms.\\
    High internal validity.
  \end{keybox}
\end{column}
\begin{column}{0.48\textwidth}
  \begin{warningbox}[Weakness]
    Do students in a lab game generalise to workers in firms?
  \end{warningbox}
\end{column}
\end{columns}
\end{column}
\begin{column}{0.42\textwidth}
  \vspace{0.5em}
  \begin{examplebox}[Classic examples]
    \small
    \begin{itemize}
      \item Ultimatum games (fairness)
      \item Public goods experiments
      \item Auction theory tests
      \item Cognitive bias elicitation
    \end{itemize}
  \end{examplebox}
\end{column}
\end{columns}
\end{frame}

\begin{frame}[t]{Field Experiments: \citeauthor{Banerjee2014_health} (\citeyear{Banerjee2014_health})}
\begin{examplebox}[Bundling Health Insurance and Microfinance in India]
  \small
  \begin{itemize}
    \item \textbf{Setting:} Microfinance borrowers in rural India
    \item \textbf{Treatment:} Randomly offered bundled health insurance at a subsidised price
    \item \textbf{Outcome:} Take-up rates for health insurance
    \item \textbf{Puzzle:} Even heavily subsidised insurance had very low take-up
  \end{itemize}
\end{examplebox}
\vspace{0.4em}
\begin{keybox}[The main finding]
There cannot be \emph{adverse selection} if there is no \emph{demand}.\\[0.3em]
Low demand is not driven by risk type — people simply do not value the product
enough to take it up, even for free.
\end{keybox}
\source{\citet{Banerjee2014_health}}
\end{frame}

\begin{frame}[t]{Field Experiments: Strengths and Challenges}
\begin{columns}[T]
\begin{column}{0.52\textwidth}
\begin{itemize}
  \setlength\itemsep{0.4em}
  \item[\textcolor{okabegreen}{\textbf{+}}] Real subjects, real decisions, real stakes
  \item[\textcolor{okabegreen}{\textbf{+}}] Results are directly policy-relevant
  \item[\textcolor{okabegreen}{\textbf{+}}] Randomisation is still under researcher control
  \item[\textcolor{okabevermillion}{\textbf{--}}] Expensive and slow to run
  \item[\textcolor{okabevermillion}{\textbf{--}}] Contamination: treated and control units may interact
  \item[\textcolor{okabevermillion}{\textbf{--}}] Attrition: participants drop out
  \item[\textcolor{okabevermillion}{\textbf{--}}] Ethical constraints on randomisation
\end{itemize}
\end{column}
\begin{column}{0.45\textwidth}
  \begin{warningbox}[Attrition example]
    \small
    We randomise a job training programme. After 6 months, 20\% of the
    \emph{treated} group has dropped out --- but only the least motivated did so.
    Comparing survivors to controls over-estimates the programme's benefit.
  \end{warningbox}
\end{column}
\end{columns}
\end{frame}

\begin{frame}[t]{Survey Experiments: What Makes Them Experiments?}
\begin{itemize}
  \setlength\itemsep{0.5em}
  \item A survey experiment randomly varies \textbf{what respondents are shown or told}
        within the survey instrument itself.
  \item The survey is the randomisation device.
  \item Common designs:
    \begin{itemize}
      \item \textbf{Information treatment:} Randomly show some respondents a fact or statistic;
            measure how it shifts beliefs or behaviour.
      \item \textbf{Vignette design:} Randomly vary characteristics of a hypothetical
            scenario; measure how responses change.
      \item \textbf{List experiment:} Measure sensitive attitudes without direct questioning.
    \end{itemize}
\end{itemize}
\vspace{0.3em}
\begin{intuitionbox}
The key is that randomisation happens \emph{inside} the survey — across respondents
who were recruited into the same sample. Standard causal inference logic applies.
\end{intuitionbox}
\end{frame}

\begin{frame}[t]{Designing a Survey Experiment: \citeauthor{Stantcheva2023_surveys} (\citeyear{Stantcheva2023_surveys})}
\begin{examplebox}[How to Run Surveys]
  \small
  \begin{itemize}
    \item Surveys can create \emph{identifying variation} that does not exist in observational data.
    \item Key design principle: the randomised element must be \emph{orthogonal}
          to all other respondent characteristics.
    \item Common pitfalls:
      \begin{itemize}
        \item Demand effects: respondents guess the ``right'' answer.
        \item Inattention: not all respondents read carefully.
        \item Hypothetical bias: stated preferences $\neq$ revealed preferences.
      \end{itemize}
    \item Solutions: attention checks, incentive-compatible designs, pre-registration.
  \end{itemize}
\end{examplebox}
\source{\citet{Stantcheva2023_surveys}}
\end{frame}

\begin{frame}[t]{Survey Experiments: Strengths and Cautions}
\begin{columns}[T]
\begin{column}{0.52\textwidth}
\begin{itemize}
  \setlength\itemsep{0.4em}
  \item[\textcolor{okabegreen}{\textbf{+}}] Study beliefs, attitudes, and preferences at large scale
  \item[\textcolor{okabegreen}{\textbf{+}}] Very fast and inexpensive
  \item[\textcolor{okabegreen}{\textbf{+}}] Can investigate topics hard to randomise in the field
        (e.g., policy beliefs, perceptions of inequality)
  \item[\textcolor{okabevermillion}{\textbf{--}}] Hypothetical bias: stated $\neq$ revealed preferences
  \item[\textcolor{okabevermillion}{\textbf{--}}] Demand effects and social desirability bias
  \item[\textcolor{okabevermillion}{\textbf{--}}] Sample representativeness concerns
\end{itemize}
\end{column}
\begin{column}{0.45\textwidth}
  \begin{intuitionbox}
    \small
    Best use of survey experiments: studying \textbf{beliefs and perceptions}
    where the respondent's subjective state \emph{is} the object of interest —
    not just their future behaviour.
  \end{intuitionbox}
\end{column}
\end{columns}
\end{frame}

\begin{frame}[t]{Preview: \citeauthor{Settele2022_gender} (\citeyear{Settele2022_gender}) — Session D1B}
\begin{examplebox}[How Do Beliefs about the Gender Wage Gap Affect Policy Demand?]
  \small
  \begin{itemize}
    \item \textbf{Question:} Do people who believe the gender wage gap is larger
          support more redistribution policies aimed at closing it?
    \item \textbf{Design:} Randomly inform a subset of respondents about the
          \emph{true} magnitude of the gender wage gap.
    \item \textbf{Outcome:} Demand for redistribution/policy intervention.
    \item \textbf{Key insight:} Information treatment creates exogenous variation
          in \emph{beliefs} — allowing identification of beliefs $\to$ policy demand.
  \end{itemize}
\end{examplebox}
\vspace{0.3em}
\muted{\small Session D1B (13:00): We will replicate the main tables of this paper in STATA.}
\source{\citet{Settele2022_gender}}
\end{frame}

% ============================================================
\sectionslide{5. Complications in Practice}{When things do not go as planned}
% ============================================================

\begin{frame}[t]{Non-Compliance: When Assignment $\neq$ Treatment}
\begin{itemize}
  \setlength\itemsep{0.5em}
  \item In a perfect experiment, everyone assigned to treatment ($Z_i = 1$) takes it,
        and no one in control ($Z_i = 0$) does.
  \item Reality: some assigned to treatment do not take it (\textbf{never-takers}).
        Some in control obtain it anyway (\textbf{always-takers}).
\end{itemize}
\vspace{0.3em}
\begin{columns}[T]
\begin{column}{0.48\textwidth}
  \begin{keybox}[One-sided non-compliance]
    Only the assignment $\to$ treatment direction fails.\\
    No one in control gets the treatment.\\
    (Common in field experiments.)
  \end{keybox}
\end{column}
\begin{column}{0.48\textwidth}
  \begin{warningbox}[Two-sided non-compliance]
    Both directions fail.\\
    Some treated opt out; some controls obtain treatment.\\
    Harder to handle.
  \end{warningbox}
\end{column}
\end{columns}
\end{frame}

\begin{frame}[t]{Intent-to-Treat (ITT)}
\begin{itemize}
  \setlength\itemsep{0.5em}
  \item The \textbf{Intent-to-Treat} estimate compares outcomes by \emph{assignment},
        ignoring actual take-up:
  \[
    \ITT = \Econd{Y_i}{Z_i=1} - \Econd{Y_i}{Z_i=0}
  \]
  \item This is always \textbf{identified} — $Z_i$ was randomly assigned.
  \item ITT is the relevant estimate for a policy question like:
    \emph{``What happens if we send everyone an invitation to the programme?''}
\end{itemize}
\vspace{0.3em}
\begin{intuitionbox}
ITT is conservative: it dilutes the treatment effect by including
non-compliers. If compliance is low, ITT can substantially understate
the effect on those who actually take up.
\end{intuitionbox}
\end{frame}

\begin{frame}[t]{LATE: The Local Average Treatment Effect}
\begin{itemize}
  \setlength\itemsep{0.5em}
  \item The \textbf{compliers} are the units who take treatment if and only if assigned
        ($D_i = 1$ when $Z_i = 1$, and $D_i = 0$ when $Z_i = 0$).
  \item The \textbf{Local Average Treatment Effect} is the ATE for compliers:
  \[
    \LATE = \frac{\ITT}{\Econd{D_i}{Z_i=1} - \Econd{D_i}{Z_i=0}}
          = \frac{\text{reduced form}}{\text{first stage}}
  \]
  \item The denominator is the share of the population who are compliers.
\end{itemize}
\vspace{0.3em}
\begin{keybox}[LATE identifies the effect for a specific subgroup]
We learn the treatment effect \emph{only for compliers} — those whose treatment
status is actually changed by the assignment. LATE $\neq$ ATE unless all
units are compliers.
\end{keybox}
\end{frame}

\begin{frame}[t]{Who Are the Compliers?}
\begin{itemize}
  \setlength\itemsep{0.5em}
  \item We \emph{cannot} identify compliers at the individual level.
  \item We can describe their \textbf{characteristics on average}
        (complier characteristic analysis).
  \item Potential outcomes for four types of units:
\end{itemize}
\vspace{0.3em}
\begin{center}
\footnotesize
\begin{tabular}{lcc}
\toprule
\textbf{Type} & \textbf{$D_i$ if $Z_i=1$} & \textbf{$D_i$ if $Z_i=0$} \\
\midrule
Complier & 1 & 0 \\
Always-taker & 1 & 1 \\
Never-taker & 0 & 0 \\
Defier & 0 & 1 \\
\bottomrule
\end{tabular}
\end{center}
\vspace{0.4em}
\begin{warningbox}[Monotonicity assumption]
LATE requires \textbf{no defiers} (monotonicity): treatment assignment either
increases or has no effect on take-up, but never decreases it.
\end{warningbox}
\end{frame}

\begin{frame}[t]{Attrition}
\begin{itemize}
  \setlength\itemsep{0.5em}
  \item \textbf{Attrition} occurs when units leave the study before the endline.
  \item \textbf{Random attrition} (equally likely in treated and control):
    \begin{itemize}
      \item Reduces power (smaller effective sample) but does not introduce bias.
    \end{itemize}
  \item \textbf{Differential attrition} (attrition rates differ by treatment):
    \begin{itemize}
      \item Survivors in each group are \emph{selected} differently.
      \item Breaks the comparability randomisation created.
      \item Can bias estimates in either direction.
    \end{itemize}
\end{itemize}
\vspace{0.3em}
\begin{keybox}[Lee Bounds]
When differential attrition is present, we can bound the treatment effect
by trimming the distribution of outcomes to make the two groups comparable.
\citet{AngristPischke2009} provide the estimator.
\end{keybox}
\end{frame}

\begin{frame}[t]{Spillovers and SUTVA Violations in Field Settings}
\begin{itemize}
  \setlength\itemsep{0.5em}
  \item When treated and control units \textbf{interact}, the independence that
        randomisation creates is broken.
\end{itemize}
\vspace{0.3em}
\begin{warningbox}[Example: Jessoe \& Rapson (2014)]
  Suppose treated households with in-home monitors \emph{tell their neighbours}
  about high-price days. Control households then also reduce consumption.
  The estimated treatment effect would be \emph{attenuated} — we would
  under-estimate the true information effect.
\end{warningbox}
\vspace{0.3em}
\begin{columns}[T]
\begin{column}{0.5\textwidth}
  \begin{keybox}[Solutions]
    \small
    \begin{itemize}
      \item Randomise at a higher level (village vs.\ household)
      \item Geographic buffer zones between treated/control
      \item Explicitly model spillovers
    \end{itemize}
  \end{keybox}
\end{column}
\begin{column}{0.47\textwidth}
  \begin{intuitionbox}
    \small
    If spillovers are a concern, the appropriate unit of randomisation
    is the \emph{community}, not the individual.
  \end{intuitionbox}
\end{column}
\end{columns}
\end{frame}

\begin{frame}[t]{The Randomisation Check (Balance Test)}
\begin{itemize}
  \setlength\itemsep{0.5em}
  \item Always check balance on \textbf{pre-treatment covariates}.
  \item Standard approach: regress each covariate $X_k$ on $D_i$ and test significance.
    \begin{itemize}
      \item Report a balance table (``Table 1'') in every empirical paper.
    \end{itemize}
  \item In a large sample, we expect $\sim5\%$ of tests to reject by chance.
  \item An \emph{F-test} on joint significance is more informative than individual $t$-tests.
\end{itemize}
\vspace{0.3em}
\begin{warningbox}[What balance failure means]
If the treatment group differs significantly from the control group on
pre-treatment covariates, something went wrong with randomisation —
or with data collection. \textbf{Investigate before proceeding.}
\end{warningbox}
\end{frame}

\begin{frame}[t]{Hawthorne Effects and Demand Effects}
\begin{itemize}
  \setlength\itemsep{0.5em}
  \item \textbf{Hawthorne effect:} Participants change behaviour simply because
        they know they are being observed, regardless of treatment.
    \begin{itemize}
      \item Leads to over-estimation of treatment effects if both groups are
            aware they are in a study.
    \end{itemize}
  \item \textbf{Demand effects:} Respondents guess the researcher's hypothesis
        and behave accordingly.
    \begin{itemize}
      \item Especially problematic in lab experiments and survey experiments.
    \end{itemize}
\end{itemize}
\vspace{0.3em}
\begin{intuitionbox}[Mitigation strategies]
  \small
  \begin{itemize}
    \item Use blind or double-blind protocols where feasible
    \item Include placebo conditions
    \item Use administrative data rather than self-reports
    \item Pre-register hypotheses (reduces demand effects in surveys)
  \end{itemize}
\end{intuitionbox}
\end{frame}

% ============================================================
\sectionslide{6. Validity}{What experiments can and cannot tell us}
% ============================================================

\begin{frame}[t]{Internal Validity}
\begin{itemize}
  \setlength\itemsep{0.5em}
  \item An experiment has \textbf{internal validity} if the estimated treatment effect
        is a valid causal estimate \emph{within the study's sample and context}.
  \item Threats to internal validity:
    \begin{itemize}
      \item Randomisation failure (e.g., non-random attrition, imperfect randomisation)
      \item Non-compliance (if not properly handled)
      \item Spillovers / SUTVA violations
      \item Contamination of the control group
      \item Measurement error in the outcome
    \end{itemize}
\end{itemize}
\vspace{0.3em}
\begin{keybox}
A study with high internal validity lets us say: ``In this study,
treatment $D$ caused a change of $X$ units in outcome $Y$ for the sample.''
That is already a strong claim.
\end{keybox}
\end{frame}

\begin{frame}[t]{External Validity}
\begin{itemize}
  \setlength\itemsep{0.5em}
  \item An experiment has \textbf{external validity} if its findings generalise
        beyond the study's specific sample and context.
  \item Common external validity concerns:
    \begin{itemize}
      \item \textbf{LATE vs. ATE:} We estimated the effect for compliers —
            are they representative of the population of interest?
      \item \textbf{Sample representativeness:} Lab with students
            $\neq$ population of workers.
      \item \textbf{Context:} Effect in rural India may not hold in urban UK.
      \item \textbf{Scale:} Small-scale pilot may not generalise when rolled out
            at scale (general equilibrium effects).
    \end{itemize}
\end{itemize}
\vspace{0.3em}
\begin{warningbox}
High internal validity $\not\Rightarrow$ high external validity.
A perfectly run experiment may answer a narrowly defined question
that does not map cleanly to the policy we care about.
\end{warningbox}
\end{frame}

\begin{frame}[t]{The LATE Problem: For Whom?}
\begin{itemize}
  \setlength\itemsep{0.5em}
  \item Experiments typically identify the \LATE\ — the effect for compliers.
  \item Compliers are \emph{not} a random sample of the population.
  \item Sometimes this is \textbf{exactly the right group}:
    \begin{itemize}
      \item A job training programme targets people who would not seek it
            without encouragement — compliers are the policy-relevant group.
    \end{itemize}
  \item Sometimes it is \textbf{not the group we care about}:
    \begin{itemize}
      \item If only motivated, low-risk individuals comply, LATE is optimistic
            relative to the effect on the full population.
    \end{itemize}
\end{itemize}
\vspace{0.3em}
\begin{intuitionbox}
Always ask: ``Whose treatment effect am I identifying?''
The answer shapes how to interpret results for policy.
\end{intuitionbox}
\end{frame}

\begin{frame}[t]{Aggregating Evidence}
\begin{itemize}
  \setlength\itemsep{0.5em}
  \item A single experiment — however well-designed — is rarely definitive.
  \item External validity improves through \textbf{replication}:
    \begin{itemize}
      \item Same design in different settings (geography, time period, population)
      \item Similar design studying a related question
    \end{itemize}
  \item \textbf{Meta-analysis} and \textbf{systematic reviews} aggregate results
        across studies, weighting by precision.
  \item Divergent findings across replications are informative:
    \begin{itemize}
      \item They identify moderating factors — when and where does the effect hold?
    \end{itemize}
\end{itemize}
\vspace{0.3em}
\muted{\small Pre-registration (AEA RCT Registry, OSF) strengthens credibility by preventing selective reporting.}
\end{frame}

\begin{frame}[t]{The Experimental Benchmark}
\begin{keybox}[The right way to think about experiments]
  Experiments are not \emph{always} feasible, ethical, or sufficient.
  But they set the \textbf{benchmark} against which we evaluate
  quasi-experimental designs.\\[0.5em]
  When an experiment is impossible, the question becomes:
  \emph{``What features of my design approximate what a random experiment would give me?''}
\end{keybox}
\vspace{0.5em}
\begin{center}
\small
\begin{tabular}{lcc}
\toprule
& \textbf{Internal validity} & \textbf{External validity} \\
\midrule
Well-run RCT & \textcolor{okabegreen}{\textbf{High}} & Depends on design \\
Field experiment & High & Better than lab \\
Survey experiment & High & Depends on sample \\
Quasi-experiment & Depends & Often high \\
\bottomrule
\end{tabular}
\end{center}
\end{frame}

% ============================================================
\sectionslide{Summary and Preview}{}
% ============================================================

\begin{frame}[t]{Key Concepts: What You Should Be Able to Do}
\begin{itemize}
  \setlength\itemsep{0.45em}
  \item Define and distinguish $\Yi{0}$, $\Yi{1}$, $\tau_i$, $\ATE$, $\ATT$
  \item Explain why the \textbf{fundamental problem of causal inference} prevents
        direct observation of $\tau_i$
  \item Explain how \textbf{randomisation} eliminates selection bias
  \item State \textbf{SUTVA} and give an example of a violation
  \item Distinguish \textbf{ITT}, \textbf{LATE}, and the role of \textbf{compliance}
  \item Identify threats to \textbf{internal} and \textbf{external} validity
        in a given experimental design
  \item Classify a study as lab, field, or survey experiment and evaluate trade-offs
\end{itemize}
\end{frame}

\begin{frame}[t]{The Causal Inference Toolkit: Days 2 and 3}
\vspace{0.5em}
\begin{itemize}
  \setlength\itemsep{0.5em}
  \item Experiments are powerful but often \textbf{infeasible}:
    \begin{itemize}
      \item Ethical constraints (we cannot randomise poverty)
      \item Cost and time
      \item Policy decisions already made
    \end{itemize}
  \item \textbf{Quasi-experimental methods} exploit naturally occurring variation
        to approximate a randomised experiment.
\end{itemize}
\vspace{0.4em}
\begin{center}
\small
\begin{tabular}{ll}
\toprule
\textbf{Day 2} & Regression Discontinuity and Difference-in-Differences \\
\textbf{Day 3} & Instrumental Variables and Propensity Score Methods \\
\bottomrule
\end{tabular}
\end{center}
\vspace{0.4em}
\begin{intuitionbox}
All quasi-experimental methods ask the same question as an RCT:
\emph{``Can I construct a valid comparison group for the treated units?''}
\end{intuitionbox}
\end{frame}

\begin{frame}[t]{Session D1B: Replicating \citeauthor{Settele2022_gender} (\citeyear{Settele2022_gender})}
\vspace{0.3em}
\textbf{What we will do in STATA:}
\vspace{0.3em}
\begin{itemize}
  \setlength\itemsep{0.4em}
  \item Load and explore the dataset (survey structure, key variables)
  \item Verify balance: are the information-treatment and control groups comparable
        on pre-treatment characteristics?
  \item Estimate the main OLS regression:
    does the information treatment shift policy preferences?
  \item Check robustness: different subgroups, alternative outcomes
\end{itemize}
\vspace{0.4em}
\begin{keybox}[Think about this before the lab]
  The information treatment is the ``instrument'' for beliefs.\\
  How would you check that the treatment actually changed beliefs
  (the first stage), and that belief changes then shift policy demand?
\end{keybox}
\source{\citet{Settele2022_gender}}
\end{frame}

\begin{frame}[plain, noframenumbering]
\vfill
\begin{center}
  {\large\bfseries\color{bathnav}Questions?}\\[1em]
  {\small\color{nearblack!70}Session D1B begins at 13:00.}\\[0.5em]
  {\scriptsize\color{nearblack!50}
    Please open STATA and load the dataset before we reconvene.}
\end{center}
\vfill
\end{frame}

% ============================================================
% BIBLIOGRAPHY
% ============================================================

\begin{frame}[allowframebreaks, noframenumbering]{References}
  \footnotesize
  \bibliography{../Bibliography_base}
\end{frame}

% ============================================================
\end{document}
% ============================================================
